\documentclass{article}
\usepackage{icomp2026_conference,times}

\usepackage[T1]{fontenc}
\usepackage[utf8]{inputenc}
\usepackage[english]{babel}
\usepackage{hyperref}
\usepackage{url}
\usepackage{booktabs}
\usepackage{nicefrac}
\usepackage{microtype}
\usepackage{xcolor}

%%%

\usepackage{subcaption}
\usepackage{graphicx}
\usepackage{multirow}
\usepackage{amsmath,amssymb,amsfonts}
\usepackage{amsthm}
\usepackage{mathrsfs}
\usepackage{textcomp}
\usepackage{manyfoot}
\usepackage{algorithm}
\usepackage{algorithmicx}
\usepackage{algpseudocode}
\usepackage{listings}

\newtheorem{theorem}{Theorem} % continuous numbers
%%\newtheorem{theorem}{Theorem}[section] % sectionwise numbers
%% optional argument [theorem] produces theorem numbering sequence instead of independent numbers for Proposition
\newtheorem{proposition}[theorem]{Proposition}% 
\newtheorem{lemma}{Lemma}% 
%%\newtheorem{proposition}{Proposition} % to get separate numbers for theorem and proposition etc.

\newtheorem{example}{Example}
\newtheorem{remark}{Remark}

\newtheorem{definition}{Definition}
\newtheorem{criterion}{Criterion}
\newtheorem{assumption}{Assumption}
\algnewcommand\Input{\item[\textbf{Input:}]}
\algnewcommand\Output{\item[\textbf{Output:}]}


\usepackage[capitalize, russian, english]{cleveref}

\crefname{definition}{definition}{definitions}
\Crefname{definition}{Definition}{Definitions}
\crefname{algorithm}{algorithm}{algorithms}
\Crefname{algorithm}{Algorithm}{Algorithms}
\crefname{section}{section}{sections}
\Crefname{section}{Section}{Sections}
\crefname{appendix}{appendix}{appendix}
\Crefname{appendix}{Appendix}{Appendix}
\crefname{fig}{fig}{fig}
\Crefname{fig}{Fig}{Fig}

%%%

\input{math_commands.tex}


\title{Predicting Grokking from 1D Training Logs via Intrinsic Attractor Dimension Estimation}


% The \author macro works with any number of authors. There are two commands
% used to separate the names and addresses of multiple authors: \And and \AND.
%
% Using \And between authors leaves it to LaTeX to determine where to break the
% lines. Using \AND forces a line break at that point. So, if LaTeX puts 3 of 4
% authors names on the first line, and the last on the second line, try using
% \AND instead of \And before the third author name.


\author{
  Nickolay Karlov\\
  MIPT\\
  Moscow, Russia\\
  \texttt{karlov.na@phystech.edu}\\
  \And
  Alexey Kravatskiy\\
  MIRIAI\\
  Moscow, Russia\\
  \texttt{kravatskii.a@miriai.org}\\
}

\icomparxivcopy

\begin{document}


\maketitle

\begin{abstract}
    We analyze delayed generalization (grokking) through the lens of empirical dynamic modeling. We hypothesize that the process of neural network training can be viewed as the evolution of a forced dynamical system possessing an attractor, which corresponds to a local minimum of the loss function, and that grokking represents a structural phase transition to a low-dimensional attractor. We apply the Levina-Bickel maximum likelihood estimator (MLE) of intrinsic dimension to 1D scalar logs of the weight norm, analyzing the resulting dimensionality trajectories to enable the early prediction of grokking. Results on modular addition and composition tasks over the non-abelian  $S_5$ symmetric group demonstrate that grokking is characterized by an initial increase in attractor dimensionality followed by a sustained reduction, which supports the hypothesis that a topological collapse underlies delayed generalization.
\end{abstract}

\textbf{Keywords:} grokking, Levina-Bickel method, phase space, dynamical systems, neural network optimization.


\section{Introduction}\label{sec:intro}

In classical statistical learning theory, the central paradigm is the bias-variance tradeoff. For a long time, it was believed that if, during the optimization process, the empirical risk (training error) drops to zero while the generalization gap remains large or begins to grow, the model irreversibly overfits, which requires stopping the training. However, in 2022, the phenomenon of grokking \citep{power2022grokking} was discovered—an effect of delayed generalization where a model, having been in a state of overfitting for thousands of iterations, undergoes a phase transition and achieves 100\% accuracy on the test set. This phenomenon challenges classical approaches to training monitoring and highlights the need for new methods for early prediction of generalization ability.

In modern literature, several approaches to studying grokking have emerged. The first direction, based on mechanistic interpretability \citep{nanda2023progress, chughtai2023toy}, investigates grokking as a process of forming structured algebraic representations (e.g., Fourier bases) within the weight matrices. Although this "white-box" approach provides a deep understanding of the mechanics, it is computationally unscalable and requires the creation of unique progress measures for each new task. The second direction focuses on the macroscopic dynamics of the loss landscape ("black-box" approach). Recent works \citep{liu2022omnigrok, thilak2022slingshot} analyze the influence of the weight norm and the slingshot effect, while the authors of \citep{notsawo2023predicting} propose using spectral signatures of the loss (Hjorth parameters) to predict the phenomenon.

However, existing macroscopic methods possess a significant limitation: they only capture statistical changes in the system's metrics, yet they remain uninterpretable in terms of the internal complexity of the model.

In this paper, we propose a new statistical framework relying on the Levina-Bickel maximum likelihood estimator (MLE) of intrinsic dimension for the early prediction and structural description of grokking. In the rigorous formulation of stochastic optimization, the task consists of minimizing the expected risk $\E_{\xi \sim \gD}[\Ls(\vw, \xi)]$, where the random variable $\xi_t$ corresponds to the selection of a specific mini-batch at iteration $t$.

Because modern adaptive optimizers (such as SGD with momentum or AdamW) utilize the history of previous gradients, the weight trajectory $\vw_t \in \R^D$ does not intrinsically form a Markov process. Therefore, we formalize the learning process as a non-autonomous dynamical system in an extended phase space:
\begin{equation}
    \vs_{t+1} = \Phi(\vs_t, \xi_t), \quad \text{where } \vs_t = (\vw_t, \vh_t) \in \sS.
\end{equation}
Here, $\vs_t$ is the full state of the system at step $t$, comprising both the model parameters $\vw_t$ and the vector of internal optimizer variables $\vh_t$ (e.g., the exponential moving averages of the first and second moments in AdamW). The deterministic operator $\Phi$ defines the state update rule, and the sequence $\{\xi_t\}$ acts as the external stochastic forcing. It is precisely this gradient noise $\xi_t$ that causes the trajectory to deviate from the ideal full-batch gradient descent, forcing the optimizer to continuously explore the local geometry of the loss landscape.

To reconstruct the properties of this system from scalar logs $y_t = \phi(\vs_t) \in \R$, we rely on embedding theorems. The presence of hidden momentum variables $\vh_t$ only increases the dimensionality of the original space $\sS$, but does not prevent reconstruction if the system converges to a low-dimensional attractor. To account for the external influence (learning rate scheduler) and the stochasticity of the forcing $\xi_t$, the classical basis of Takens' embedding theorem \citep{takens1981detecting} is extended by Stark's theorems on \textit{delay embeddings} for systems with stochastic forcing \citep{stark1999delay, stark2003delay}. Stark proved that for such systems, the delay map preserves diffeomorphism, which provides us with the mathematical justification to estimate the effective dimensionality of the optimizer's attractor exclusively from one-dimensional macroscopic logs.

We consider grokking as a structural topological transition. In the memorization phase, the optimizer's trajectory constitutes a high-dimensional stochastic random walk, during which the system allocates redundant degrees of freedom to memorize the training set rather than to discover generalized rules. In the generalization phase, $L_2$ regularization (weight decay) forces the system to discard noisy degrees of freedom and build a compact algebraic structure within the weights. Topologically, this is equivalent to the compression of the optimizer's attractor into a low-dimensional submanifold.

To verify the proposed hypothesis and analyze the phenomenon of grokking, we conduct a series of computational experiments. As the primary mathematical tool, we apply the Levina-Bickel maximum likelihood estimation (MLE) of intrinsic dimension \citep{levina2004maximum}. Training is performed on algebraic datasets using a 1-layer Transformer, as proposed in the Omnigrok study \citep{liu2022omnigrok}. In the first experiment, the problem of modular addition (a commutative group) is solved; in the second, we tackle the composition operation within the non-abelian permutation group $S_5$. Such a choice of test environments allows us to evaluate the proposed method both for its ability to provide early prediction of delayed generalization and for its sensitivity to changes in the hidden algebraic complexity of the task.






\section{Problem Statement}\label{sec:problem}

Consider a neural network model $\gM$ parameterized by weights $\vw \in \R^D$. In the presence of $L_2$ regularization, the model is optimized to minimize the regularized expected risk:
\begin{equation}
    \gJ(\vw) = \E_{\xi \sim \gD}[\Ls(\vw, \xi)] + \frac{\lambda}{2} \|\vw\|_2^2,
\end{equation}
where $\Ls(\vw, \xi)$ is the primary loss function (e.g., cross-entropy), $\xi$ denotes the stochastic mini-batch selection, and $\lambda \ge 0$ is the weight decay (WD) coefficient. During optimization, the regularization term induces an explicit shrinkage force $-\eta \lambda \vw_t$ in the parameter update step (where $\eta$ is the learning rate). As we demonstrate in our experiments, this restoring force plays a critical role in grokking by restricting the unbounded growth of the parameters and inducing a topological collapse of the system's degrees of freedom.

Because modern adaptive optimizers (e.g., SGD with momentum or AdamW) maintain internal state histories, the weight trajectory $\vw_t$ alone does not form a Markov process. Therefore, let $\vs_t = (\vw_t, \vh_t) \in \sS$ be the extended state of the system at step $t \in \{1, \dots, T\}$, where $\vh_t$ represents the optimizer's internal variables. The evolution of this system is governed by the recurrence rule:
\begin{equation}
    \vs_{t+1} = \Phi(\vs_t, \xi_t),
\end{equation}
where $\Phi$ is the deterministic optimizer step derived from the gradient $\nabla \gJ(\vw_t)$, and $\xi_t$ acts as the external stochastic forcing induced by the random mini-batch selection and hyperparameter scheduling.

In practice, the full high-dimensional state $\vs_t$ is inaccessible for computationally efficient analysis. Instead, we observe a set of one-dimensional scalar time series $\vm_t =[m_1(t), \dots, m_K(t)]^\top$ generated during optimization. Each macroscopic metric (e.g., training loss, weight norm, or gradient norm) acts as an observation function $m_i(t) = \phi_i(\vs_t, \xi_t) \in \R$. In this regard, the \textbf{global research objective} is to reconstruct the topological properties of the hidden attractor of $\vs_t$ and to develop an early predictor of phase transitions during optimization, relying exclusively on the 1D observation series $m_i(t)$.

The primary target phase transition in our work is the phenomenon of delayed generalization. Let us introduce its formal description.

\begin{definition}[Grokking]\label{def:grokking}
Let $\text{Acc}_{train}(\vw_t)$ and $\text{Acc}_{val}(\vw_t)$ denote the model's accuracy on the training and validation sets, respectively. Due to stochastic optimizer noise, instantaneous metric values are prone to random fluctuations (e.g., spurious accuracy spikes early in training). Let $\overline{\text{Acc}}(t)$ denote the smoothed metric (a moving average over a local time window). For a given small error threshold $\epsilon > 0$ (set to $0.05$ in this work), we define two critical optimization moments:
\begin{itemize}
    \item \textbf{Memorization moment ($t_{mem}$):} the minimum step $t$ starting from which the network robustly interpolates the training data:
    $$t_{mem} = \min \left\{ t \mid \forall \tau \ge t, \; \overline{\text{Acc}}_{train}(\vw_\tau) \ge 1 - \epsilon \right\}.$$
    Typically, at this moment, the validation accuracy remains at the level of random guessing ($\overline{\text{Acc}}_{val}(\vw_{t_{mem}}) \approx 0$).
    
    \item \textbf{Generalization moment ($t_{gen}$):} the minimum step $t$ starting from which a stable increase in generalization ability is recorded:
    $$t_{gen} = \min \left\{ t \mid \forall \tau \ge t, \; \overline{\text{Acc}}_{val}(\vw_\tau) \ge 1 - \epsilon \right\}.$$
\end{itemize}
The optimization process is classified as grokking if there is an anomalous delay between these events: $t_{gen} \gg t_{mem}$.
\end{definition}

As \cref{def:grokking} implies, on the interval $t \in (t_{mem}, t_{gen})$, classical macroscopic accuracy metrics reside on a quasi-stationary plateau. During this period, the optimizer outwardly appears trapped in an overfitting regime; however, hidden structural reorganizations occur within the model weights. We hypothesize that the transition from memorization to generalization constitutes a topological collapse of the optimizer's degrees of freedom.

To detect this collapse, we introduce $\hat{d}(t)$—an estimate of the effective (intrinsic) dimensionality of the reconstructed attractor at time $t$. Based on this, we formulate a formal criterion for successful early prediction.

\begin{definition}[Grokking prediction criterion]\label{def:criterion}
A successful early warning signal of grokking is characterized by the formation of a stable downward trend in the effective dimensionality $\hat{d}(t)$, culminating in its stabilization at a new, lower-dimensional level strictly before the actual increase in validation accuracy:
\begin{equation}
    \exists t_{drop} \in (t_{mem}, t_{gen}) \quad \text{such that} \quad \forall t \in [t_{drop}, t_{gen}]: \; \; \hat{d}(t) \ll \hat{d}(t_{mem}).
\end{equation}
Furthermore, on the interval $(t_{mem}, t_{drop})$, a macroscopically monotonic decline of $\hat{d}(t)$ without significant upward rebounds must be observed.
\end{definition}

Meeting the conditions of \cref{def:criterion} in practice implies that the topological dimensionality collapse indeed precedes generalization, and that the scalar observation functions $m_i(t)$ contain sufficient information for its early detection.









\section{Theoretical Justification and Proposed Method}\label{sec:method}

To solve the stated problem, we propose a framework relying on the theory of nonlinear dynamical systems. This section justifies the theoretical possibility of extracting the geometry of a high-dimensional system from a 1D scalar log and presents the algorithm for estimating the effective dimensionality $\hat{d}(t)$.

\subsection{Learning as a Forced Dynamical System (Embedding Theorems)}

In nonlinear dynamics theory, an \textbf{attractor} $\mathcal{A} \subset \sW$ is understood as a compact invariant subset of the phase space to which the system's trajectories asymptotically converge. In the context of machine learning, the gradient optimization process seeks to minimize the loss function. In this regard, we rely on the assumption that over time, the weight trajectory $\vw_t$ loses redundant degrees of freedom and converges to a certain local attractor (e.g., a stationary stochastic cloud around a flat minimum). Let the true dimensionality of this attractor be $m$, where $m \ll D$.

According to Takens' embedding theorem \citep{takens1981detecting}, for autonomous systems, the topology of the entire high-dimensional attractor $\mathcal{A}$ can be reconstructed from a single scalar observation series $X(t) \in \{m_i(t)\}$ using delay embeddings:
\begin{equation}
    \vx_t =[X(t), X(t-\tau), \dots, X(t-(E-1)\tau)]^\top \in \R^E,
\end{equation}
provided that the dimensionality of the embedding space $E \ge 2m + 1$.

To formalize the algorithms, we introduce the $\text{DelayEmbedding}$ construction operator, which maps the original one-dimensional time series $\{X(t)\}_{t=1}^N$ to a sequence of vectors in the reconstructed phase space:
\begin{equation}
    \text{DelayEmbedding}(\{X(t)\}_{t=1}^N, E, \tau) = \{ \vx_t \}_{t=(E-1)\tau + 1}^{N},
\end{equation}
where each vector $\vx_t \in \R^E$ is constructed according to Eq. (3). The result of this operator is the reconstructed phase manifold $M \subset \R^E$, consisting of $|M| = N - (E-1)\tau$ points.

A fundamental mathematical requirement of embedding theorems is the \textit{genericity} of the observation function $\phi(\cdot)$. From the perspective of differential topology, this means that the metric function must not possess degenerate symmetries that glue distinct points of the phase space together during projection. For example, considering the extended state $\vs_t = (\vw_t, \vh_t)$, the aggregated weight norm $||\vw_t||_2$ possesses spherical symmetry, measuring only the radial distance of the parameters from the origin while entirely ignoring the optimizer's history $\vh_t$. In certain regimes, it can act as a "degenerate observable," losing information about the algebraic structure (angular coordinates) of the system. Conversely, nonlinear metrics directly dependent on the logit distribution (such as the loss function $\Ls$) act as generic observables that preserve the high-dimensional topology of the extended system.

However, gradient descent in deep learning is not an autonomous system. The influence of random mini-batch selection, combined with hyperparameter schedules (learning rate schedulers), formally violates the classical Takens conditions. To mathematically justify our approach in this non-autonomous setting, we rely on Stark's embedding theorems for forced dynamical systems with mixed forcing \citep{stark1999delay, stark2003delay}.

In modern optimization theory, the extended learning process $\vs_{t+1} = \Phi(\vs_t, \omega_t, \xi_t)$ almost always occurs under the influence of two types of forcing. The deterministic forcing $\omega_t$ (e.g., a learning rate schedule) obeys Stark's theorem for deterministic systems \citep{stark1999delay}, which guarantees the preservation of topology under a smooth external driver. In turn, the stochastic forcing $\xi_t$ (mini-batch gradient noise) is approximated by a multivariate Gaussian process \citep{mandt2017stochastic}, and the selection of a batch from a finite set is described in terms of Iterated Function Systems \citep{stark2003delay}. For both cases, Stark and co-authors proved that the delay map preserves diffeomorphism properties. This fundamentally justifies the validity of measuring the dimensionality of the hidden attractor in the extended state space $\sS$ by operating exclusively on scalar logs in the reconstructed phase space $\R^E$, regardless of the presence of internal optimizer memory, deterministic schedules, or stochastic noise.

\subsection{Effective Dimensionality Estimation of the Phase Space}\label{sec:theory_dim}

To detect phase transitions, we need to dynamically estimate the attractor dimension of the system. In our framework, we reviewed both classical and modern methods, ranking them by their robustness to the specifics of neural network optimization (mini-batch noise):

\paragraph{False Nearest Neighbors (FNN) \citep{kennel1992determining} and Cao's Method \citep{cao1997practical}.}
These are geometric "bottom-up" approaches. The essence of FNN lies in the sequential "unfolding" of the attractor. Given an insufficient embedding dimension, distant segments of the high-dimensional trajectory are projected close to each other, forming "false" neighbors. The method measures the distance between a point and its nearest spatial neighbor, and checks how much they diverge when transitioning to dimension $d+1$. Cao's method is a rigorous mathematical development of this approach that analyzes the averaged ratio of distances, eliminating the need for manual threshold selection.

However, a fundamental limitation of both geometric methods is the strict requirement of \textit{global recurrence} of the trajectory. To identify the complex shape of the attractor, the system must regularly return to the vicinity of its past states. During neural network optimization, we frequently deal with transient processes—for example, when the loss function drops monotonically. In the absence of recurrence, the nearest spatial neighbor of a point $\vx_t$ becomes its own temporal neighbor $\vx_{t-1}$. Because such points do not "scatter" in any dimensionality, geometric methods degenerate, erroneously perceiving the high-dimensional descent trajectory as a simple 1D line, yielding a trivial estimate of $\hat{d}=1$.

\paragraph{Simplex Projection Maximization \citep{sugihara1990nonlinear}.}
The optimal $E$ is sought as the argument that maximizes the autoregressive correlation $\rho$ when predicting the future state $y_{t+1}$ based on its neighbors in $\R^E$. This topologically guarantees the best unfolding of the attractor without overfitting to noise; however, its discrete (integer-valued) resolution coarsens the detection of smooth phase transitions.

\paragraph{Maximum Likelihood Estimation (MLE Intrinsic Dimension) \citep{levina2004maximum}.}
Selected as the primary analytical tool in our framework. In contrast to geometric "bottom-up" methods, the Levina-Bickel method projects the time series directly into a redundant high-dimensional space (e.g., $E=15$) and estimates the local intrinsic dimension $\hat{d}_{x_t}$ statistically, based on the Poisson distribution of distances to the $k$ nearest neighbors. Due to space constraints, the full pseudocode is deferred to \cref{app:alg}. Briefly, for each embedded point within a sliding window $W$, the algorithm computes a local dimensionality estimate using the log-ratio of distances to its $k$ nearest neighbors, and averages these local estimates to output the effective dimensionality $\hat{d}(t)$.

This method returns a \textit{continuous} value that is computed locally, does not require global recurrence of the trajectory, and demonstrates phenomenal robustness to microscopic mini-batch noise.

\subsection{Grokking as a Topological Collapse}\label{sec:theory_grokking}

Applying the described mathematical apparatus, we formalize grokking as a structural phase transition in the loss landscape.
\begin{itemize}
    \item \textbf{Memorization phase:} The network "memorizes" the training data as uncorrelated noise. The optimizer trajectory $\vw_t$ wanders along a flat, high-dimensional manifold (high-entropy state). During this period, the effective dimensionality $\hat{d}(t)$ computed by the MLE algorithm will be maximal ($\hat{d} \gg 1$).
    \item \textbf{Generalization phase (Grokking):} Under the influence of $L_2$ regularization (weight decay), the network discards noisy features and "crystallizes" the optimal algebraic algorithm within the structure of its weights.
\end{itemize}

Finding a generalizing solution strictly restricts the degrees of freedom of the system. Geometrically, this means that the dynamical system's attractor contracts into a low-dimensional submanifold. Consequently, tracking the effective dimensionality $\hat{d}(t)$ over a sliding window should act as an early indicator of the system's topological simplification, without requiring the analysis of heavy Hessian matrices.










\section{Experiments}\label{sec:exp}

To test the hypothesis that grokking is a structural phase transition accompanied by a reduction in attractor dimensionality, and to validate our statistical framework as an early predictor of grokking, we conducted a series of computational experiments.

The primary experiments presented in this section utilize the 1-layer Omnigrok Transformer architecture \citep{liu2022omnigrok}, which was specifically designed to facilitate a smoother and more distinct grokking effect. Training for these cases was performed using the AdamW optimizer with mini-batch stochasticity. As a standard practice in deep learning, this stochastic noise increases the variance of the scalar metric series, allowing phase reconstruction algorithms to more accurately estimate the true dimensionality of the phase space. In contrast, our initial baseline experiments, which were conducted using a standard Transformer architecture and full-batch gradient descent (Full-Batch GD), revealed the problem of a "degenerate observable" due to the loss dropping to zero and the resulting lack of variance; these results and their analysis are deferred to \cref{app:exp}.

In addition to the baseline metrics ($\Ls$, accuracy), the algorithm tracks the gradient norm $\|g\|_2$, the weight norm $\|\vw\|_2$, and the cosine similarity between gradients at consecutive iterations. For the phase space reconstruction, the time delay $\tau$ is typically set to a fixed constant (e.g., $\tau=1$ or $\tau=10$), while the effective embedding dimension $E$ is computed dynamically using the Maximum Likelihood Estimation (MLE) method.

\subsection{Modular Addition}\label{sec:exp_mod}

In the first experiment, we tackled the modular addition task (modulo $p=113$). The weight norm ($||\vw||_2$) was chosen as the target metric for the MLE method, as it preserves rich variance throughout the entire memorization phase and, equally importantly, remains completely independent of the validation set.

\cref{fig:grokking_flat} presents the results of the experiment. In the top row (with $L_2$ regularization enabled, $WD=1.0$), a perfect topological collapse is observed: the effective dimensionality $\hat{d}(t)$ collapses synchronously with the weight norm long before the Accuracy metric registers generalization. In the bottom row (without regularization, $WD=0.0$), grokking does not occur, the network enters perpetual memorization, and the MLE method correctly shows a persistently high dimensionality of the phase space. This proves the method's robustness against false positives.

In summary, our framework successfully predicts grokking by capturing the topological collapse of the attractor dimension. However, it is evident that the weight norm is a non-generic (symmetric) observable: while it clearly reflects the dimensionality drop without relying on the validation set, it fails to convey the true algebraic dimension of the attractor. Furthermore, our predictor on this specific task does not surpass the early-warning accuracy of a predictor based directly on the weight norm and $\Ls_{val}$, which is attributed to the simplicity of the modular addition task itself.

\begin{figure}[htbp]
    \centering
    % Top row: Grokking (WD=1.0)
    \begin{subfigure}[b]{0.32\textwidth}
        \includegraphics[width=\textwidth]{images/mod_wd1_acc.png}
        \caption{WD=1: Dim \& Accuracy}
        \label{fig:grokking_flat_a}
    \end{subfigure}
    \hfill
    \begin{subfigure}[b]{0.32\textwidth}
        \includegraphics[width=\textwidth]{images/mod_wd1_loss.png}
        \caption{WD=1: Dim \& Loss}
        \label{fig:grokking_flat_b}
    \end{subfigure}
    \hfill
    \begin{subfigure}[b]{0.32\textwidth}
        \includegraphics[width=\textwidth]{images/mod_wd1_norm.png}
        \caption{WD=1: Dim \& Norm}
        \label{fig:grokking_flat_c}
    \end{subfigure}
    
    \vspace{0.3cm}
    % Bottom row: No grokking (WD=0.0)
    \begin{subfigure}[b]{0.32\textwidth}
        \includegraphics[width=\textwidth]{images/mod_wd0_acc.png}
        \caption{WD=0: Dim \& Accuracy}
        \label{fig:grokking_flat_d}
    \end{subfigure}
    \hfill
    \begin{subfigure}[b]{0.32\textwidth}
        \includegraphics[width=\textwidth]{images/mod_wd0_loss.png}
        \caption{WD=0: Dim \& Loss}
        \label{fig:grokking_flat_e}
    \end{subfigure}
    \hfill
    \begin{subfigure}[b]{0.32\textwidth}
        \includegraphics[width=\textwidth]{images/mod_wd0_norm.png}
        \caption{WD=0: Dim \& Norm}
        \label{fig:grokking_flat_f}
    \end{subfigure}
    \caption{Analysis of grokking on modular addition in the presence of stochastic noise. The plots compare the dynamics of dimensionality $E$ (derived from the weight norm series) against Accuracy, Loss, and Weight Norm in regimes with generalization (top row) and without it (bottom row).}
    \label{fig:grokking_flat}
\end{figure}

\subsection{Non-Abelian Groups and Complex Algebraic Representations ($S_5$)}\label{sec:exp_s5}

To verify the universality of our method, we escalated the task complexity to the composition of permutations within the non-abelian $S_5$ symmetric group. From the perspective of mechanistic interpretability, neural networks generalize on finite group tasks by learning their linear irreducible representations \citep{chughtai2023toy}.

While commutative modular addition is projected by the network into 2D circles (where the minimal representation dimension is 2), the structure of the non-abelian group $S_5$ is significantly more complex. According to algebraic representation theory, the dimensions of the irreducible representations of $S_5$ are 1, 1, 4, 4, 5, 5, and 6. Since 1D representations lose information about the permutations themselves (i.e., they are unfaithful), the network must utilize weight subspaces of dimension at least 4 to flawlessly compute the composition of any elements. 

In the context of nonlinear dynamics, this means the true attractor dimension after grokking should collapse but plateau at a level no lower than 4 components. This provides a rigorous test for our method's ability to "blindly" measure the algebraic complexity of the task.

As seen in \cref{fig:grokking_s5}, the proposed method successfully tackles this challenge. During the chaotic search phase, the dimensionality $E$ remains at a high level. Once the network forms a generalizing algebraic rule, the dimensionality predictably collapses. The control experiment without Weight Decay (bottom row) reaffirms that the dimensionality does not drop when grokking is absent.

\begin{figure}[htbp]
    \centering
    % Top row: Grokking S5 (WD=1.0)
    \begin{subfigure}[b]{0.32\textwidth}
        \includegraphics[width=\textwidth]{images/s5_wd1_acc.png}
        \caption{S5 WD=1: Dim \& Accuracy}
        \label{fig:grokking_s5_a}
    \end{subfigure}
    \hfill
    \begin{subfigure}[b]{0.32\textwidth}
        \includegraphics[width=\textwidth]{images/s5_wd1_loss.png}
        \caption{S5 WD=1: Dim \& Loss}
        \label{fig:grokking_s5_b}
    \end{subfigure}
    \hfill
    \begin{subfigure}[b]{0.32\textwidth}
        \includegraphics[width=\textwidth]{images/s5_wd1_norm.png}
        \caption{S5 WD=1: Dim \& Norm}
        \label{fig:grokking_s5_c}
    \end{subfigure}
    
    \vspace{0.3cm}
    % Bottom row: No grokking S5 (WD=0.0)
    \begin{subfigure}[b]{0.32\textwidth}
        \includegraphics[width=\textwidth]{images/s5_wd0_acc.png}
        \caption{S5 WD=0: Dim \& Accuracy}
        \label{fig:grokking_s5_d}
    \end{subfigure}
    \hfill
    \begin{subfigure}[b]{0.32\textwidth}
        \includegraphics[width=\textwidth]{images/s5_wd0_loss.png}
        \caption{S5 WD=0: Dim \& Loss}
        \label{fig:grokking_s5_e}
    \end{subfigure}
    \hfill
    \begin{subfigure}[b]{0.32\textwidth}
        \includegraphics[width=\textwidth]{images/s5_wd0_norm.png}
        \caption{S5 WD=0: Dim \& Norm}
        \label{fig:grokking_s5_f}
    \end{subfigure}
    \caption{Analysis of composition in the $S_5$ group. The dimensionality collapse reliably captures the structural reorganization of the network, even for the complex multidimensional topology of a non-abelian group.}
    \label{fig:grokking_s5}
\end{figure}

Of particular interest is the dimensionality estimation using the validation loss function (\cref{fig:grokking_val_loss_s5}). This proves that the observed topological collapse is independent of the chosen metric and is not a subjective artifact. By evaluating a generic observable ($\Ls_{val}$), we can correctly measure the number of active components in the model. For $S_5$, the dimensionality collapses to $E \approx 4$, which perfectly aligns with the minimal required representation dimension. Conversely, for $WD=0$, the dimensionality not only fails to drop but actively increases, physically reflecting the escalating chaos of eternal memorization.

\begin{figure}[htbp]
    \centering
    \begin{subfigure}[b]{0.49\textwidth}
        \includegraphics[width=\textwidth]{images/s5_wd1_val_loss.pdf}
        \caption{S5 WD=1: Dim from validation loss}
        \label{fig:grokking_val_loss_s5_a}
    \end{subfigure}
    \hfill
    \begin{subfigure}[b]{0.49\textwidth}
        \includegraphics[width=\textwidth]{images/s5_wd0_val_loss.pdf}
        \caption{S5 WD=0: Dim from validation loss}
        \label{fig:grokking_val_loss_s5_b}
    \end{subfigure}
    \vspace{0.3cm}
    \caption{Analysis of the $S_5$ group. The dimensionality $E$ reconstructed from the validation loss successfully estimates the true algebraic complexity of the algorithm (plateauing at $\hat{d} \approx 4$).}
    \label{fig:grokking_val_loss_s5}
\end{figure}

\subsection{Theoretical Interpretation of Results}\label{sec:exp_theory}

An analysis of the obtained plots allows us to draw the following theoretical conclusions.

\paragraph{The nature of grokking.} The observed dimensionality collapse (\cref{fig:grokking_s5_a,fig:grokking_val_loss_s5_a})  demonstrates that grokking is not merely \textbf{quantitative loss minimization}, but a qualitative structural reconfiguration of the dynamical system: a transition from a high-entropy "memorization" attractor to a low-dimensional "algebraic rule" attractor.

\paragraph{Optimizer metrics predictability.} The experiments demonstrated that delayed generalization can be diagnosed relying exclusively on internal network metrics (such as weight norm) \cref{fig:grokking_flat_a,fig:grokking_s5_a}. This suggests that the transition to generalization is an \textbf{intrinsic topological property} of the optimization trajectory and is not merely an artifact of a specific held-out validation set.

\paragraph{The role of genericity of the observable.} Evaluating the $WD=0$ regime revealed the critical importance of choosing a generic observable. Weight norm, a symmetric observable, remains low during overfitting, as the model simply finds a trivial solution and scales up the weights, keeping $E \approx 1$ \cref{fig:grokking_flat_d,fig:grokking_s5_d}. However, a generic observable ($\Ls_{val}$) correctly captures the high dimensionality of chaotic overfitting and accurately projects the algebraic complexity of the learned algorithm (e.g., $E \ge 4$ for $S_5$) (\cref{fig:grokking_val_loss_s5_a}).

\section{Discussion}\label{sec:disc}

The results of our computational experiments confirm the fundamental hypothesis of this work: by framing deep neural network optimization as a nonlinear dynamical system, one can predict and qualitatively describe grokking through tracking the attractor dimension. However, transferring phase reconstruction methods to machine learning has revealed specific limitations and physical effects that require separate consideration.

\paragraph{The Problem of the Degenerate Observable and the Role of Stochasticity.}
During the experiments with full-batch gradient descent (Full-Batch GD), we encountered difficulties in estimating the attractor dimension: all the 1D scalar metrics we used degenerated into a constant. From the perspective of nonlinear dynamics, ideal noise-free optimization collapses the trajectory into a rectilinear ray, depriving methods (such as MLE) of the variance necessary to estimate the landscape curvature. Thus, when analyzing dimensionality, one must ensure that the series does not degenerate. This can be mitigated by subsampling the series, though this may negatively impact the estimation accuracy. Therefore, our method performs best in systems where the observation series do not degenerate, i.e., in systems with inherent stochasticity. The high-frequency noise of SGD forces the optimizer to continuously explore the local geometry of the attractor. Furthermore, experiments have shown that the weight norm ($||\vw||_2$) is a much more informative and stable "observable" for the latent dynamics than the training loss, which is highly prone to degeneration.

\paragraph{Topology Dependence on the Choice of Observable (Observability).}
One of the most important empirical results of our work is demonstrating the dependence of the reconstructed dimensionality $E$ on the choice of the scalar metric (observation function). In dynamical systems theory, Takens' embedding theorem requires the observation function to be "generic" and free of degenerate symmetries.

Our experiments have demonstrated that the weight norm ($||\vw||_2$), being a symmetric and thus non-generic observable, distorts the absolute value of the true dimensionality. The dynamics of the norm are primarily determined by the balance of two global forces: the error gradient (which increases the weights) and $L_2$ regularization (which shrinks them). Because of this, the effective dimensionality of the $||\vw||_2$ series always collapses to low values ($E \approx 2\dots 4$), regardless of the complexity of the solved task, as this metric is "blind" to the complex angular structure of the representations.

In contrast, the loss function ($\Ls_{val}$ or the noisy $\Ls_{train}$) depends directly on the logits and is sensitive to the internal algebraic structure of the network. Works on the direct spectral analysis of weights \citep{nanda2023progress} show that for modular addition, the network forms 2D circles (Fourier bases). MLE estimation on the loss series confirms this topologically: the trajectory collapse halts at $E \approx 2$. In turn, for the non-abelian group $S_5$, whose irreducible representations require spaces of higher dimensionality \citep{chughtai2023toy}, the dimensionality collapse, measured specifically on the loss series, predictably halts at a higher level of $E \approx 4\dots 6$. Thus, we have proven that for a "blind" estimation of the algebraic complexity of the learned algorithm, it is necessary to use metrics directly related to the distribution of the network outputs, or to aggregate parameter statistics more carefully. It is also worth noting that it is preferable to use specifically $\Ls_{val}$, since $\Ls_{train}$ degenerates rapidly.

\paragraph{Independence from the Validation Set and the Number of Active Components.}
From a practical standpoint, monitoring the validation loss ($\Ls_{val}$) remains a faster and cheaper predictor of the onset of grokking, at least in the simple experiments conducted. However, macroscopic metrics are merely symptoms of improved generalization ability. In turn, the collapse of the effective dimensionality $E$ provides rigorous mathematical evidence of a structural phase transition—a process in which the network truncates noisy degrees of freedom and compresses the trajectory into a compact submanifold. Furthermore, our dimensionality-based predictor fundamentally does not require a validation set. 
% This observation paves the way for fundamentally new methods of network pruning and model compression (e.g., rank selection for LoRA), based on the true number of degrees of freedom of the trained network, computed dynamically.



\section{Conclusion}\label{sec:concl}

In this work, we proposed and verified a novel statistical method for the early diagnosis of grokking, based on the topological reconstruction of the neural network’s phase space. By treating stochastic gradient descent as a forced dynamical system, we empirically demonstrated that the phenomenon of delayed generalization constitutes a structural phase transition to a low-dimensional state. Tracking this process using the Levina-Bickel estimation (MLE Intrinsic Dimension) allows for capturing the collapse of the effective dimensionality of the attractor long before the jump in validation accuracy, thus acting as a reliable early predictor of generalization.

The proposed approach is characterized by high universality and practical applicability. It is shown that the topological collapse can be reliably detected from the internal scalar logs of the optimizer (e.g., the weight norm), which enables monitoring the generalization capability without the need for a validation set. Furthermore, experiments on the tasks of modular addition and composition in the non-abelian group $S_5$ confirmed that the residual attractor dimension strictly aligns with the theoretical algebraic complexity of the learned algorithm when using a generic observable ($E \approx 2$ and $E \approx 4 \dots 6$, respectively). Since the method operates on a "black-box" principle and relies exclusively on 1D logs, it can be easily transferred to arbitrary architectures and tasks without adaptation.

The primary conclusion of this study is that one-dimensional scalar series contain sufficient information for diagnosing hidden phase transitions during the training process. This paves the way for the development of computationally cheap methods for the monitoring and structural optimization of deep models without the need for resource-intensive analysis of hidden representations or Hessian matrices.


%%%%%%%%%%%%%%%%%%%%%%%%%%%%%%%%%%%%%%%%%%%%%%%%%%%%%%%%%%%%

\bibliographystyle{icomp2026_conference}
\bibliography{references}

\appendix
\section{MLE Intrinsic Dimension Estimation Algorithm}\label{app:alg}

The full pseudocode for the Levina-Bickel maximum likelihood estimation (MLE) of intrinsic dimension, as utilized in \cref{sec:theory_dim}, is provided in \cref{alg:mle_id}.

\begin{algorithm}[htbp]
\caption{MLE Intrinsic Dimension Estimation (Levina-Bickel)}\label{alg:mle_id}
\begin{algorithmic}[1]
\Input Observation time series $\gD_{obs}$, delay $\tau$, maximum embedding dimension $E_{max}$, window $W$, neighbors $k$
\Output Effective dimensionality $\hat{d}(t)$
\State $M \gets \text{DelayEmbedding}(\gD_{obs}, E_{max}, \tau)$ \Comment{Embedding into $\R^{E_{max}}$}
\For{each point $\vx_t \in M$ within the sliding window $W$}
    \State Find $k$ nearest neighbors with distances $r_1 \le r_2 \le \dots \le r_k$
    \State $\hat{d}_{x_t} \gets \left[ \frac{1}{k-1} \sum_{j=1}^{k-1} \log \left( \frac{r_k}{r_j} \right) \right]^{-1}$ \Comment{Local estimation via radius $r_k$}
\EndFor
\State $\hat{d}(t) \gets \frac{1}{|W|} \sum_{\vx_t \in W} \hat{d}_{x_t}$ \Comment{Averaging dimensionality over the window $W$}
\State \Return $\hat{d}(t)$
\end{algorithmic}
\end{algorithm}
\section{Baseline Experiment with Full-Batch Gradient Descent}\label{app:exp}

This appendix presents the results of initial experiments conducted on the modular addition task using full-batch gradient descent (Full-Batch GD). This experiment served as a baseline to confirm the initial hypothesis.

As seen in \cref{fig:grokking_full_batch}, in this setting, we were able to capture the drop in the effective dimensionality (MLE ID) of the $\Ls_{train}$ series immediately preceding the jump in validation accuracy (from $E \approx 4-5$ to $E \approx 2.8$).

However, this experiment revealed an important methodological problem. Under full-batch optimization, the correlation between the phase space dimension $E$ and the grokking effect became less evident. From a theoretical perspective, this is explained by the "degenerate observable" problem: in the absence of the stochastic noise of mini-batches, the optimizer's trajectory degenerates into a deterministic smooth curve, and the training loss drops almost to zero, depriving the MLE method of the variance necessary to estimate the true dimension of the space $d^*(t)$.

This very observation motivated the transition to stochastic optimization (mini-batch) and the use of a non-generic observable (the weight norm), the results of which are presented in the main \cref{sec:exp}.

\begin{figure}[H]
    \centering
    \begin{subfigure}[b]{0.42\textwidth}
    \centering
    \includegraphics[width=\textwidth]{images/grokking_accuracy.png}
    \caption{Dynamics of Train/Val Accuracy}
    \label{fig:grokking_full_batch_a}
    \end{subfigure} \hfill
    \begin{subfigure}[b]{0.42\textwidth}
    \centering
    \includegraphics[width=\textwidth]{images/grokking_dimension.png}
    \caption{Dimensionality Collapse (MLE ID)}
    \label{fig:grokking_full_batch_b}
    \end{subfigure}
\caption{The grokking phenomenon in the baseline experiment (Full-Batch). The transition from memorization to generalization is accompanied by a decrease in the topological complexity of the $\Ls_{train}$ trajectory; however, the estimation suffers from a lack of variance due to the degeneration of the series.}
\label{fig:grokking_full_batch}
\end{figure}

\end{document}